\section{Introduction}

The architecture of psychiatric assessment is undergoing a fundamental shift from the episodic, clinic-bound snapshot to the continuous, ecologically valid stream of behavior in situ. This paradigm, formalized as \textit{digital phenotyping} \cite{onnela2016harnessing}, promises to bridge the ``quantification gap'' that has historically hindered mental health care—the inability to observe the patient during the $99.9\%$ of their life spent outside the clinician's office. By leveraging the ubiquity of smartphones and wearable sensors, researchers have sought to create a ``Digital Smoke Detector'' \cite{nock2026predicting} capable of identifying the subtle precursors to relapse, self-harm, and suicidal ideation before they manifest in catastrophic clinical outcomes. However, a decade after the field's inception, a critical systemic failure has emerged, threatening the validity of these predictive models: the phenomenon of \textit{Informative Missingness} (Missing Not At Random, MNAR), which we formalize here through the lens of the ``Digital Vacuum Paradox.''

\subsection{The Dawn of Digital Phenotyping and the Promise of Continuous Sensing (2015–2018)}
The early years of digital phenotyping were characterized by a tech-utopian optimism. Following the seminal work of Onnela and colleagues \cite{onnela2016harnessing}, the field transitioned from reliance on retrospective self-reports—long known to be plagued by recall bias and mood-congruent memory—to the passive collection of high-frequency behavioral data. Mobility patterns (GPS), social interactions (anonymized call and text logs), and physical activity (accelerometry) were hailed as the new ``digital biomarkers'' of the psychiatric phenotype. Insel \cite{insel2017digital} famously argued that digital phenotyping represented the first true ``microscope'' for behavioral health, allowing for the observation of intra-individual variability with millisecond precision.

During this period, studies such as those by Barnett et al. \cite{barnett2018relapse} demonstrated the feasibility of predicting relapse in schizophrenia by identifying deviations from a patient's personalized behavioral baseline. These models typically employed Gaussian processes or recurrent neural networks (RNNs) to smooth the inherent noise in sensor data. The underlying assumption, often unstated, was that as long as the sensors remained active, the model would remain vigilant. The ``smoke detector'' was assumed to be always-on, awaiting the first flicker of behavioral anomaly.

\subsection{The Great Dropout: The Crisis of Attrition and Engagement (2019–2023)}
As digital phenotyping moved from small-scale feasibility pilots to large-scale longitudinal cohorts, the reality of the ``Great Dropout'' began to manifest. By 2021, a meta-analysis of mHealth studies in psychiatry revealed a staggering attrition rate: over $50\%$ of participants in naturalistic studies ceased engagement within the first 30 days, a cliff that Torous et al. \cite{torous2021growing} described as the primary barrier to the clinical translation of digital health. Unlike the attrition seen in traditional clinical trials, which is often driven by study-specific factors, dropout in digital phenotyping is intrinsically linked to the underlying psychiatric pathology.

This era marked a shift from viewing missing data as a technical nuisance to recognizing it as a fundamental clinical signal. Chen et al. \cite{chen2022attrition} noted that patients with the highest severity of depressive symptoms were precisely those most likely to stop charging their devices, uninstall monitoring apps, or ignore ecological momentary assessments (EMA). This created a selection bias where the models were trained predominantly on the ``healthy'' periods of a patient's trajectory, leaving them blind to the very crises they were designed to detect. The ``Great Dropout'' was not merely a loss of data; it was a loss of the most \textit{relevant} data.

\subsection{The 'Silent Smoke Detector' and the Informative Power of Absence}
By 2026, the work of Nock et al. \cite{nock2026predicting} provided a critical re-evaluation of the smoke detector metaphor. They argued that the most dangerous failure mode for a psychiatric monitoring system is not the false alarm, but the ``Silent Smoke Detector.'' In the context of suicide risk, the transition from high-intensity ideation to acute action is frequently accompanied by a complete cessation of digital interaction—a withdrawal from social networks, a freezing of mobility patterns, and a total silence in device telemetry.

Standard anomaly detection algorithms, which rely on observing deviations in signal (e.g., a sudden increase in night-time activity), are fundamentally incapable of detecting a crisis characterized by the \textit{absence} of signal. If the model assumes that the data-generating process is independent of the latent clinical state (Missing At Random, MAR), it will interpret a flatline in activity as a continuation of the previous state or, worse, as a sign of behavioral stability. This realization necessitated a formalization of \textit{Informative Missingness} (MNAR), where the mask of missingness $M_t$ is itself a descendant of the latent psychiatric state $S_t$ in a causal M-graph.

\subsection{OS-Level Biases and the Technological Shadow}
The challenge of attrition is further compounded by the evolving landscape of mobile operating systems. Modern smartphones employ aggressive battery optimization techniques—such as Android's ``Doze'' mode and Apple's background execution limits—to preserve hardware longevity. While beneficial for the average user, these features introduce a systematic bias into digital phenotyping. Apps that do not receive frequent user interaction are deprioritized by the OS, leading to the silent termination of background sensor collection.

This introduces a secondary layer of MNAR: the ``Technological Shadow.'' Patients who are disengaged due to low mood are punished by the OS, which kills the background processes of the monitoring app precisely because the patient is not interacting with it. Furthermore, this creates a socioeconomic bias; patients with older, lower-end hardware (characterized by poor battery health and aggressive OS-level killing of background processes) are disproportionately represented in the missing data pool. Schultebraucks et al. \cite{schultebraucks2025llms} highlighted that failing to account for these OS-level biases risks misclassifying socioeconomic hardship or technical friction as clinical deterioration, leading to inequitable care.

\subsection{Formalization of the Digital Vacuum Paradox}
We define the \textit{Digital Vacuum Paradox} as the structural inverse relationship between a patient's clinical risk and the availability of their digital evidence. Formally, let $S_t$ be the latent psychiatric state and $Y_t$ be the observed behavioral stream. The paradox posits that:
\begin{equation}
    P(M_t = 1 | S_t) \propto \text{Risk}(S_t)
\end{equation}
As $S_t$ approaches a critical threshold (e.g., the manifold of crisis or relapse), the probability of missingness $M_t$ approaches unity. This creates a ``Digital Vacuum''—a period of maximum clinical uncertainty that coincides with the period of maximum clinical need.

The paradox reveals a fundamental flaw in current SOTA imputation methods like BRITS \cite{cao2018brits} or GAIN \cite{yoon2018gain}, which attempt to fill gaps by interpolating from local temporal anchors. In psychiatric crises, the vacuum is often so prolonged that the temporal anchors are themselves lost. We refer to this as the \textit{Persistence Gap}: a threshold where the autocorrelation of the behavioral signal is broken by the duration of the disengagement. To recover the signal from the vacuum, one must move beyond interpolation to \textit{generative recovery} using a causal model of why the data is missing.

\subsection{Bridging the Gap: Clinical Data Loss Prevention (DLP)}
This paper introduces the Clinical Data Loss Prevention (DLP) Framework, a novel approach to psychiatric monitoring that treats disengagement not as a nuisance, but as a primary clinical outcome. Borrowing from enterprise IT, where DLP focuses on preventing the loss of sensitive data, our framework focuses on preventing the loss of \textit{causal signal} during the Digital Vacuum.

By integrating Latent Stochastic Differential Equations (SDEs) with a Transformer-based memory architecture, we model the latent psychiatric state as a Jump-Diffusion process. This allows us to capture the ``catastrophic'' state transitions that characterize psychiatric relapse, even when those transitions occur within the vacuum of missing data. We utilize the ``Silent Smoke Detector'' as a training objective, forcing the model to learn that prolonged absence is a strong prior for clinical risk. Through the use of synthetic clinical cohorts and causal estimands for N-of-1 data \cite{cai2026causal}, we demonstrate that it is possible to maintain a high-fidelity ``Clinical Vigilance Index'' (CVI) even when the sensors go silent. The goal is a resilient system that remains vigilant not just when the patient is using their phone, but especially when they are not.
